\section{中线和角平分线模型}

\begin{theorem}[中线模型]
    处理中线问题时，通常对中线分割出的两个小三角形分别使用余弦定理，利用互补角（如 $\angle ADB + \angle ADC = 180^\circ$）的余弦值互为相反数来建立联系。当然，还有一个重要的结论：
\end{theorem}

\begin{theorem}[阿波罗尼奥斯定理]
    和中线长有关。$b^2+c^2 = 2(m_a^2 + (\frac{a}{2})^2)$，其中 $m_a$ 是边 $a$ 上的中线。
\end{theorem}

\begin{example}
    在 $\triangle ABC$ 中，$AB=4$, $AC=6$, $BC$ 边上的中线 $AD= \sqrt{10}$，求 $BC$ 的长度。
\end{example}

\begin{tikzpicture}[scale=0.5]
    \coordinate (B) at (0,0);
    \coordinate (C) at (8,0);

    \path[name path=circle_from_B, overlay] (B) circle (4);
    \path[name path=circle_from_C, overlay] (C) circle (6);

    \path[name intersections={of=circle_from_B and circle_from_C, by=A}];

    \coordinate (D) at ($(B)!0.5!(C)$);

    \draw[thick] (B) -- (C) -- (A) node[midway, above, sloped] {6} -- (B) node[midway, above, sloped] {4};
    \draw[thick, dashed, blue] (A) -- (D) node[midway, above right, sloped] {$\sqrt{10}$};

    \node[below left] at (B) {$B$};
    \node[below right] at (C) {$C$};
    \node[above] at (A) {$A$};
    \node[below, blue] at (D) {$D$};

    \fill (A) circle (2pt);
    \fill (B) circle (2pt);
    \fill (C) circle (2pt);
    \fill (D) circle (2pt);
\end{tikzpicture}

\begin{solution}
    设 $D$ 是 $BC$ 的中点，$\angle ADB = \theta$，则 $\angle ADC = 180^\circ - \theta$。根据 $ BD = DC$，分别对两个子三角形列余弦定理即可求解。下面给出用阿波罗尼奥斯定理的证明过程：
\end{solution}

\v{2em}

\begin{proof}
    在 $\triangle ABD$ 中，由余弦定理：
    $$
        AB^2 = AD^2 + BD^2 - 2 AD \cdot BD \cos\theta
    $$
    在 $\triangle ACD$ 中，由余弦定理：
    $$
        AC^2 = AD^2 + CD^2 - 2 AD \cdot CD \cos(180^\circ - \theta)
    $$
    因为 $D$ 是中点，所以 $BD=CD$。又因为 $\cos(180^\circ - \theta) = -\cos\theta$，所以第二个式子变为：
    $$
        AC^2 = AD^2 + BD^2 + 2 AD \cdot BD \cos\theta
    $$
    将两个式子相加，消去含 $\cos\theta$ 的项：
    $$
        AB^2 + AC^2 = 2(AD^2 + BD^2)
    $$
    代入已知数据：
    $$
        4^2 + 6^2 = 2((\sqrt{10})^2 + BD^2)
    $$
\end{proof}
\vspace{2em}

\begin{theorem}[角平分线模型]
    处理角平分线问题时，最常用的方法是面积法，即大三角形的面积等于角平分线分割出的两个小三角形面积之和。
\end{theorem}

\begin{example}
    在 $\triangle ABC$ 中，$AB=3$, $AC=2$, $\angle BAC = 120^\circ$。若 $AD$ 是 $\angle BAC$ 的角平分线，求线段 $AD$ 的长度。
\end{example}
\begin{solution}
    设 $AD$ 的长度为 $x$。
    利用面积法：$S_{\triangle ABC} = S_{\triangle ABD} + S_{\triangle ACD}$。
    根据面积公式 $S = \frac{1}{2}ab \sin C$：
    $$
        S_{\triangle ABC} = \frac{1}{2} AB \cdot AC \sin 120^\circ = \frac{1}{2} AB \cdot AD \sin 60^\circ + \frac{1}{2} AC \cdot AD \sin 60^\circ.
    $$
    所以，我们有方程：
    $$
        \frac{3\sqrt{3}}{2} = \frac{3\sqrt{3}}{4}x + \frac{\sqrt{3}}{2}x
    $$
    解得线段 $AD$ 的长度为 $\frac{6}{5}$。
\end{solution}
\vspace{2em}

\begin{problem}
$\triangle A B C$ 的内角 $A, B, C$ 的对边分别为 $a, b, c$ ，其面积为 $\frac{\sqrt{3}}{4}\left(c^2-a^2-b^2\right)$ ．
\begin{enumerate}
    \item （1）求角 $C$ ；
    \item （2）若 $C$ 的角平分线交 $A B$ 于点 $D$ ，且 $C D=\sqrt{3}, B D=3 A D$ ，求 $c$ 的值．
\end{enumerate}
\end{problem}
\v{12em}

\begin{problem}
在 $\triangle A B C$ 中，内角 $A, B, C$ 所对的边分别为 $a, ~ b, ~ c$ ，其外接圆的半径为 $\sqrt{3}$ ，且 $\sin ^2 A+\sin ^2 C=\sin ^2 B-\sin A \sin C$.
\begin{itemize}
    \item （1）求 $B$；
    \item （2）若 $B$ 的角平分线交 $A C$ 于点 $D, ~ B D=\frac{\sqrt{3}}{2}$ ，点 $E$ 在线段 $A C$ 上，$E C=2 E A$ ，求 $\triangle B D E$ 的面积．
\end{itemize}
\end{problem}

\v{12em}

\begin{problem}
在 $\triangle A B C$ 中，$a, b, c$ 分别是角 $A, B, C$ 的对边，$b=2 \sqrt{3},(a+c)^2-b^2=a c$ ．
\begin{enumerate}
    \item 求角 $B$ 的大小；
    \item 若 $A D$ 是 $\angle B A C$ 的内角平分线，当 $\triangle A B C$ 面积最大时，求 $A D$ 的长．
\end{enumerate}
\end{problem}

\v{10em}

\newpage

\subsection{张角定理}

\begin{theorem}[张角定理]
    在 $\triangle \mathrm{ABC}$ 中， D 是 BC 上的一点，连结 AD ，则有：
    $$
        \frac{\sin \angle B A C}{A D}=\frac{\sin \angle B A D}{A C}+\frac{\sin \angle C A D}{A B}
    $$
    \begin{center}
        \begin{tikzpicture}[scale=1]
            \coordinate (A) at (0,2);
            \coordinate (B) at (-2,0);
            \coordinate (C) at (3,0);
            \coordinate (D) at (0.5,0);

            \draw[thick] (A) -- (B) -- (C) -- cycle;
            \draw[thick, blue] (A) -- (D);

            \node[above] at (A) {$A$};
            \node[below left] at (B) {$B$};
            \node[below right] at (C) {$C$};
            \node[below] at (D) {$D$};

            \filldraw (A) circle (2pt);
            \filldraw (B) circle (2pt);
            \filldraw (C) circle (2pt);
            \filldraw (D) circle (2pt);

            % 标注角
            \pic [draw, ->, angle eccentricity=2] {angle=B--A--D};
            \pic [draw, ->, angle eccentricity=2] {angle=D--A--C};
            \pic [draw, <-, angle eccentricity=1] {angle=B--A--C};
        \end{tikzpicture}
    \end{center}
\end{theorem}

\subsection{中线衍生出的向量模型}

\begin{theorem}[中线向量]
    \begin{minipage}[c]{0.6\textwidth}
        在三角形 $\triangle ABC$ 中：
        \begin{itemize}
            \item 设 $D$ 是 $BC$ 边的中点，则 $\overrightarrow{AD} = \frac{1}{2}(\overrightarrow{AB}+\overrightarrow{AC})$
        \end{itemize}
        这一结论可以用于求解中线长度及其与其他量之间的关系。
    \end{minipage}
    \begin{minipage}[c]{0.4\textwidth}
        \begin{tikzpicture}[scale=0.8]
            \coordinate (A) at (0,3);
            \coordinate (B) at (-1.5,0);
            \coordinate (C) at (2,0);
            \coordinate (D) at ($(B)!0.5!(C)$);

            \draw[thick] (A) -- (B) -- (C) -- cycle;
            \draw[thick, red] (A) -- (D);

            \node[above] at (A) {$A$};
            \node[below left] at (B) {$B$};
            \node[below right] at (C) {$C$};
            \node[below] at (D) {$D$};

            \node[red, right] at ($(A)!0.6!(D)$) {$\overrightarrow{AD}$};

            \filldraw (A) circle (2pt);
            \filldraw (B) circle (2pt);
            \filldraw (C) circle (2pt);
            \filldraw (D) circle (2pt);
        \end{tikzpicture}
    \end{minipage}
\end{theorem}

\begin{example}
    已知 $B = \frac{\pi}{3}$，设 $D$ 为 $A C$ 的中点，$b=2$ ；求 $B D$ 的最大值．
\end{example}

\vspace{10em}

\subsection{回顾：交叉相乘成比例}

\begin{theorem}[三点共线定理]
    设 $A, B, C$ 是平面上三点，对于直线 $BC$ 上任意一点 $D$，若存在实数 $\lambda, \mu$ 满足 $\lambda + \mu = 1$，使得 $\overrightarrow{AD} = \lambda\overrightarrow{AB} + \mu\overrightarrow{AC}$，则 $A, B, C$ 三点共线。

    反之，若点 $D$ 在直线 $BC$ 上，则一定存在实数 $\lambda, \mu$ 满足 $\lambda + \mu = 1$，使得 $\overrightarrow{AD} = \lambda\overrightarrow{AB} + \mu\overrightarrow{AC}$。

    \begin{minipage}[c]{0.65\textwidth}
        特别地，当 $D$ 是 $BC$ 的中点时，有 $\overrightarrow{AD} = \frac{1}{2}\overrightarrow{AB} + \frac{1}{2}\overrightarrow{AC}$。

        这个定理提供了用向量表示共线点的一种方式，在三角形中点问题、线性组合问题中非常有用。
    \end{minipage} %
    \begin{minipage}[c]{0.35\textwidth}
        \begin{tikzpicture}[scale=0.8]
            \coordinate (A) at (0,3);
            \coordinate (B) at (-1.5,0);
            \coordinate (C) at (2,0);
            \coordinate (D) at ($(B)!0.6!(C)$);

            \draw[thick] (A) -- (B) -- (C) -- cycle;
            \draw[thick, red] (A) -- (D);

            % 绘制向量箭头
            \draw[->,thick,blue] (A) -- ($(A)!0.8!(B)$) node[midway,above left] {$\overrightarrow{AB}$};
            \draw[->,thick,blue] (A) -- ($(A)!0.8!(C)$) node[midway,above right] {$\overrightarrow{AC}$};
            \draw[->,thick,red] (A) -- ($(A)!0.7!(D)$) node[midway,right] {$\overrightarrow{AD}$};

            \node[above] at (A) {$A$};
            \node[below left] at (B) {$B$};
            \node[below right] at (C) {$C$};
            \node[below] at (D) {$D$};

            \filldraw (A) circle (2pt);
            \filldraw (B) circle (2pt);
            \filldraw (C) circle (2pt);
            \filldraw (D) circle (2pt);
        \end{tikzpicture}
    \end{minipage}
\end{theorem}

\begin{example}
    已知 $a, ~ b, ~ c$ 分别为 $\triangle A B C$ 三个内角 $A, ~ B, ~ C$ 的对边，$(\sin B-\sin C)^2=\sin ^2(B+C)-\sin B \sin C$ ，且 $\triangle A B C$ 的面积为 $2 \sqrt{3}$ ．
    \begin{enumerate}
        \item 求 $A$ ；
        \item 若点 $D$ 在边 $BC$ 上，且 $\overrightarrow{BD}=2 \overrightarrow{DC}$ ，求 $A D$ 的最小值．
    \end{enumerate}
\end{example}

\v{12em}

\begin{example}
    已知 $\triangle A B C$ 中，三个内角分别为 $A, B, C, \cos A+2 \sin \left(\frac{3 \pi}{4}+A\right) \cos \left(\frac{3 \pi}{4}+A\right)=1$ ，且 $A \neq \frac{\pi}{2}$ ．
    \begin{enumerate}
        \item 若 $B C=1$ ，求 $\triangle A B C$ 的外接圆面积；
        \item 若 $\overrightarrow{B D}=4 \overrightarrow{C D}, \angle D A B=90^{\circ}, \triangle A B D$ 的面积为 $2 \sqrt{3}$ ，求 $\triangle A B C$ 的周长．
    \end{enumerate}
\end{example}

\v{12em}
